% !TeX root = ../master.tex

\chapter{Background}
\label{ch:background}

\section{Argument Mining}
\label{sec:argument_mining}

% - Definition: automatic identification and extraction of argumentative structures from text
% - Core subtasks: argument identification (binary), component classification, relation extraction
% - This thesis focuses on argument identification: is a sentence an argument or not?
% - Benchmark landscape: 17+ datasets across domains (medical, legal, scientific, political, financial)
% - Variability in annotation guidelines and what counts as "argument" across datasets
% - The generalization problem: models trained on one dataset fail on others

\section{Shortcut Learning}
\label{sec:shortcut_learning}

% - Geirhos et al. (2020): shortcuts = decision rules that work on i.i.d. but fail on distribution shift
% - Pervasive across deep learning: texture bias in vision, lexical cues in NLP
% - Not unique to ML: humans and animals also take shortcuts (Geirhos examples)
% - The issue is not shortcuts per se, but that they break under distribution shift
% - Jakobsen et al. (2021): spurious correlations in cross-topic argument mining
% - Feger et al. (ACL 2025) key findings:
%   - Encoders show delta <= 0.02 under word removal -> rely on content words
%   - 97% of cross-dataset experiments below benchmark performance
%   - WRAP shows largest drop (delta=0.05) -> learned some argument features
%   - Joint benchmark training helps but doesn't solve the problem

\section{Encoder vs. Decoder Architectures}
\label{sec:encoder_decoder}

% - Encoders (BERT): bidirectional attention, [CLS] aggregation
%   - Representation is somewhat order-invariant (bag-of-words-like)
%   - Explains why removing function words barely matters (delta <= 0.02)
% - Decoders (Llama, GPT): causal/autoregressive attention, left-to-right
%   - Word order enforced by attention mask, not just positional embeddings
%   - "therefore", "because" shape all subsequent computation
%   - Removing discourse markers disrupts the entire chain
% - Prediction: decoders should show larger delta under manipulation
% - Sinha et al. (2021): "Unnatural Language Inference" -- evidence for encoder order invariance

\section{Large Language Models for Argument Mining}
\label{sec:llms_am}

% - Cabessa et al. (COLING 2025): fine-tuned LLMs achieve SOTA across AM subtasks
% - Cabessa & Mushtaq (2024): GPT-4 with ICL matches fine-tuned encoder baselines
% - Gorur et al. (COLING 2025): LLMs outperform RoBERTa on relation-based AM via prompting
% - Gap: none of these test perturbation robustness or cross-dataset generalization under manipulation
% - Xie et al. (2024): "Do LLMs overcome shortcut learning?" -- general LLM shortcut evaluation
% - GoLLIE (Sainz et al., 2023): annotation guidelines as structured prompts for zero-shot IE
%   - Surpassed models fine-tuned on hundreds of datasets
%   - Direct methodological precedent for context utilization in this thesis

\section{Data Contamination in LLMs}
\label{sec:contamination}

% - LLM pre-training corpora may include public benchmark datasets
% - Golchin & Surdeanu (ICLR 2024): sentence-completion contamination test
% - Threat to validity: if LLMs have memorized GAIC sentences, zero-shot results are inflated
% - Mitigation: ROUGE-L + exact match contamination test applied to all models
% - Temporal analysis: older datasets (pre-2018) vs. newer ones
